\section{Introduction}
\label{sec:intro}

Philosophers have long asked what happens to a subject when its organization
is moved. A teleporter scans a person and builds an exact replica elsewhere. A
gradual procedure replaces a brain part by part. An upload copies the
pattern and leaves the original running
\citep{Parfit1984, Chalmers2010Singularity, Cerullo2015}. The usual question is
whether the result is the same subject. That question has resisted settlement,
and we do not try to settle it here. We ask a prior question that can be
answered exactly. Suppose the organization described by a closure architecture
for consciousness \citep{TsiokosArchitecture2026} is moved from one finite
system to a different one. What, precisely, would make the result a
\emph{continuation} of the source rather than a \emph{copy} of it, and which of
the conditions involved can be checked?

We show that the question splits into three, and that the three must be
checked separately. The first is \emph{realization}: does the target system
satisfy the conditions of the architecture, on its own state space? The second
is \emph{lineage}: does the episode pass from source to target through handoffs
in which the two systems are live at the same time, with no gap? The third is
\emph{autonomy}: once the source and every aid used for the transfer are
withdrawn, does the target keep itself going? A detached copy can pass the
first and the third and fail the second. A resumed stream can match the
source at both ends and still fail the second. A target that runs only while
an outside resource is supplied fails the third. Similarity of endpoints
settles none of the three.

Our main result is a precise form of this separation. We define three checks
whose predicates are defined from the systems' own data. Realization and
lineage are computed from it. Autonomy additionally requires a supplied proof
of an invariance obligation. We prove that a certificate passing them is sound: it yields
realization on the target, a strict live-overlap lineage, and autonomy after
withdrawal (Theorem~\ref{thm:continuation}, the continuation theorem). The
proofs are elementary. Local pieces that agree where they overlap glue into a
single realization on the target. A one-step agreement between source and
target extends to equality of their observation laws over every finite sequence
of actions. A certificate of invariance
yields a strategy that keeps the target going. The work lies elsewhere: in
formulating checks strong enough to separate continuation from its
look-alikes, and in exhibiting exact finite systems that pass them and others
that fail. The theorem comes in two versions. In the operational version, each
of the three conclusions can fail while the other two hold. In the
realization-preserving version, autonomy is defined to include realization,
so it implies it. Our passing examples include a chain of handoffs through
three distinct systems. On these models four of the eight conditions hold only
in a simplified form, which we describe.

\begin{figure}[tbp]
\centering
% F1: schematic episode axes; ordinals label stages/checks, not measured data.
% Verdicts R/L/A: continuation +/+/+; detached +/x/+; gap -/x/-;
% one cell deleted x/+/+; recharge withdrawn +/+/x.
\begingroup
\newcommand{\episodeStageOne}{1}
\newcommand{\episodeStageTwo}{2}
\newcommand{\episodeStageThree}{3}
\newcommand{\episodeStageFour}{4}
\begin{tikzpicture}[x=1cm,y=1cm,font=\footnotesize,
 bar/.style={line width=1.5pt}, muted/.style={black!35},
 aid/.style={line width=.65pt}, frame/.style={draw=black!20,line width=.35pt},
 badge/.style={circle,draw,fill=white,inner sep=1.7pt,font=\scriptsize},
 smalllabel/.style={font=\scriptsize,align=center},
 loop/.style={-{Stealth[length=1.4mm]},line width=.6pt}]
\path[use as bounding box] (0,-7.60) rectangle (15.35,4.22);
\node[anchor=west,font=\scriptsize,text=black!65] at (0,4.08) {schematic};
\node at (3.8,4.00) {before};
\node at (7.5,4.00) {overlap};
\node[align=center] at (11.1,4.00) {after withdrawal};
\node at (3.8,3.55) {\episodeStageOne};
\node at (6.25,3.55) {\episodeStageTwo};
\node at (8.75,3.55) {\episodeStageThree};
\node at (11.1,3.55) {\episodeStageFour};
\node[font=\footnotesize\bfseries] at (13.0,4.0) {R};
\node[font=\footnotesize\bfseries] at (13.85,4.0) {L};
\node[font=\footnotesize\bfseries] at (14.7,4.0) {A};
\node[badge] at (13.0,3.55) {\episodeStageOne};
\node[badge] at (13.85,3.55) {\episodeStageTwo};
\node[badge] at (14.7,3.55) {\episodeStageThree};
% The same cut crosses all five rows; row frames mark their right edges.
\fill[black!3] (10,-.10) rectangle (12.1,2.55);
\fill[black!10] (5,1.05) rectangle (10,2.55);
\draw[frame] (2.65,-.10) rectangle (12.1,2.55);
\node[badge] at (5.58,3.14) {\episodeStageTwo};
\node at (7.5,3.14) {live overlap, no gap};
\node[smalllabel] at (7.5,2.80) {each step supported on both};
\draw (6.25,2.43)--(6.25,2.62);
\draw (8.75,2.43)--(8.75,2.62);
\node[anchor=east] at (2.45,2.30) {source};
\node[anchor=east] at (2.45,1.30) {target};
\draw[bar] (2.8,2.30)--(10,2.30);
\draw[bar,densely dashed,black!35] (2.8,1.30)--(5,1.30);
\node[smalllabel,text=black!65] at (3.8,1.57) {latent};
\draw[bar] (5,1.30)--(12.40,1.30);
% Badge lies on the target only: no connecting stroke to the source.
\draw[rounded corners=2pt,fill=white,line width=.5pt] (5.38,1.10) rectangle (7.83,1.50);
\node[badge] at (5.57,1.30) {\episodeStageOne};
\node[smalllabel] at (6.80,1.84) {eight conditions, on the\\target's own dynamics};
\draw[fill=white,line width=.6pt] (5,1.18)--(5.12,1.30)--(5,1.42)--(4.88,1.30)--cycle;
\node[smalllabel] at (4.85,.88) {activation / receiver};
\draw[aid] (5,.64)--(10,.64);
\node[smalllabel] at (7.75,.84) {transfer aid};
\draw[aid] (2.8,.06)--(12.40,.06);
\node[smalllabel] at (7.5,-.29) {target's own recharge};
\node[badge] at (11.78,1.80) {\episodeStageThree};
\draw[loop] (11.02,1.30) .. controls (10.73,1.72) and (11.65,1.94) .. (11.45,1.42);
\node[smalllabel] at (11.13,.85) {keeps itself\\going};
\node[font=\footnotesize\bfseries] at (13.85,.95) {continuation};
\draw[black,line width=.8pt] (9.910,2.210)--(10.090,2.390) (9.910,2.390)--(10.090,2.210);
\draw[black,line width=.8pt] (9.910,0.550)--(10.090,0.730) (9.910,0.730)--(10.090,0.550);
\fill[black] (5.95,1.3) circle (.055);
\fill[black] (6.18,1.3) circle (.055);
\fill[black] (6.41,1.3) circle (.055);
\fill[black] (6.64,1.3) circle (.055);
\fill[black] (6.87,1.3) circle (.055);
\fill[black] (7.1,1.3) circle (.055);
\fill[black] (7.33,1.3) circle (.055);
\fill[black] (7.56,1.3) circle (.055);
\draw[line width=.85pt] (12.870,1.820)--(12.965,1.725)--(13.150,1.960);
\draw[line width=.85pt] (13.720,1.820)--(13.815,1.725)--(14.000,1.960);
\draw[line width=.85pt] (14.570,1.820)--(14.665,1.725)--(14.850,1.960);
\node[anchor=west,align=left,font=\scriptsize] at (.08,-1.3) {detached\\copy};
\draw[frame] (2.65,-1.96) rectangle (12.1,-0.78);
\node[anchor=east,font=\scriptsize,text=black!50] at (2.60,-1.05) {source};
\node[anchor=east,font=\scriptsize,text=black!50] at (2.60,-1.55) {target};
\draw[aid,muted] (2.8,-1.89)--(12.4,-1.89);
\draw[bar] (2.8,-1.05)--(4.82,-1.05);
\fill (4.76,-1.11) rectangle (4.88,-0.99);
\draw[bar,muted] (5,-1.55)--(12.4,-1.55);
\draw[densely dotted,line width=.65pt] (4.82,-1.05)--(4.82,-0.71)--(12.20,-0.71)--(12.20,-1.55);
\node[smalllabel,fill=white,inner sep=1pt] at (8.3,-0.71) {same readout};
\draw[fill=white] (12.20,-1.55) circle (.055);
\draw[loop,muted] (11.02,-1.55) .. controls (10.85,-1.21) and (11.55,-1.18) .. (11.40,-1.47);
\draw[rounded corners=1pt,black!35,fill=white] (5.78,-1.69) rectangle (7.75,-1.41);
\fill[black!35] (5.95,-1.55) circle (.055);
\fill[black!35] (6.18,-1.55) circle (.055);
\fill[black!35] (6.41,-1.55) circle (.055);
\fill[black!35] (6.64,-1.55) circle (.055);
\fill[black!35] (6.87,-1.55) circle (.055);
\fill[black!35] (7.1,-1.55) circle (.055);
\fill[black!35] (7.33,-1.55) circle (.055);
\fill[black!35] (7.56,-1.55) circle (.055);
\draw[line width=.85pt] (12.870,-1.300)--(12.965,-1.395)--(13.150,-1.160);
\draw[black,line width=.8pt] (13.745,-1.405)--(13.955,-1.195) (13.745,-1.195)--(13.955,-1.405);
\draw[line width=.85pt] (14.570,-1.300)--(14.665,-1.395)--(14.850,-1.160);
\node[anchor=west,align=left,font=\scriptsize] at (.08,-2.95) {gap};
\draw[frame] (2.65,-3.61) rectangle (12.1,-2.43);
\node[anchor=east,font=\scriptsize,text=black!50] at (2.60,-2.7) {source};
\node[anchor=east,font=\scriptsize,text=black!50] at (2.60,-3.2) {target};
\draw[aid,muted] (2.8,-3.54)--(12.4,-3.54);
\fill[black!7] (7.5,-3.33) rectangle (10,-2.52);
\draw[black!35] (8.75,-2.58)--(8.75,-2.46);
\draw[bar,muted] (2.8,-2.7)--(5,-2.7) (7.5,-2.7)--(10,-2.7);
\draw[bar,muted,densely dashed] (2.8,-3.2)--(5,-3.2);
\draw[bar,muted] (7.5,-3.2)--(12.4,-3.2);
\draw[pattern=north east lines,pattern color=black!60,line width=.6pt] (5,-3.29) rectangle (7.5,-2.61);
\node[smalllabel,fill=white,inner sep=2pt] at (6.25,-2.95) {missing stage};
\draw[aid,muted] (7.5,-3.38)--(10,-3.38);
\draw[black!35,line width=.8pt] (9.910,-2.790)--(10.090,-2.610) (9.910,-2.610)--(10.090,-2.790);
\draw[black!35,line width=.8pt] (9.910,-3.470)--(10.090,-3.290) (9.910,-3.290)--(10.090,-3.470);
\draw[densely dotted,line width=.65pt] (10,-2.7)--(10.15,-2.36)--(12.20,-2.36)--(12.20,-3.2);
\node[smalllabel,fill=white,inner sep=1pt] at (11.10,-2.36) {same readout};
\draw[fill=white] (12.20,-3.2) circle (.055);
\draw[line width=.7pt] (12.88,-2.95)--(13.12,-2.95);
\draw[black,line width=.8pt] (13.745,-3.055)--(13.955,-2.845) (13.745,-2.845)--(13.955,-3.055);
\draw[line width=.7pt] (14.58,-2.95)--(14.82,-2.95);
\node[anchor=west,align=left,font=\scriptsize] at (.08,-4.6) {one cell\\deleted};
\draw[frame] (2.65,-5.26) rectangle (12.1,-4.08);
\node[anchor=east,font=\scriptsize,text=black!50] at (2.60,-4.35) {source};
\node[anchor=east,font=\scriptsize,text=black!50] at (2.60,-4.85) {target};
\fill[black!7] (5,-4.98) rectangle (10,-4.17);
\draw[black!35] (6.25,-4.23)--(6.25,-4.11);
\draw[black!35] (8.75,-4.23)--(8.75,-4.11);
\draw[aid,muted] (2.8,-5.19)--(12.4,-5.19);
\draw[bar,muted] (2.8,-4.35)--(10,-4.35) (5,-4.85)--(12.4,-4.85);
\draw[bar,muted,densely dashed] (2.8,-4.85)--(5,-4.85);
\draw[muted,fill=white,line width=.5pt] (5,-4.97)--(5.12,-4.85)--(5,-4.73)--(4.88,-4.85)--cycle;
\draw[aid,muted] (5,-5.03)--(10,-5.03);
\draw[black!35,line width=.8pt] (9.910,-4.440)--(10.090,-4.260) (9.910,-4.260)--(10.090,-4.440);
\draw[black!35,line width=.8pt] (9.910,-5.120)--(10.090,-4.940) (9.910,-4.940)--(10.090,-5.120);
\draw[rounded corners=1pt,black!35,fill=white] (5.78,-4.99) rectangle (7.75,-4.71);
\fill[black!35] (5.95,-4.85) circle (.055);
\fill[black!35] (6.18,-4.85) circle (.055);
\fill[black!35] (6.41,-4.85) circle (.055);
\draw[fill=white,dash pattern=on 1pt off .7pt,line width=.5pt] (6.64,-4.85) circle (.063);
\fill[black!35] (6.87,-4.85) circle (.055);
\fill[black!35] (7.1,-4.85) circle (.055);
\fill[black!35] (7.33,-4.85) circle (.055);
\fill[black!35] (7.56,-4.85) circle (.055);
\draw[loop,muted] (11.02,-4.85) .. controls (10.85,-4.51) and (11.55,-4.48) .. (11.40,-4.77);
\draw[black,line width=.8pt] (12.895,-4.705)--(13.105,-4.495) (12.895,-4.495)--(13.105,-4.705);
\draw[line width=.85pt] (13.720,-4.600)--(13.815,-4.695)--(14.000,-4.460);
\draw[line width=.85pt] (14.570,-4.600)--(14.665,-4.695)--(14.850,-4.460);
\node[anchor=west,align=left,font=\scriptsize] at (.08,-6.25) {recharge\\withdrawn};
\draw[frame] (2.65,-6.91) rectangle (12.1,-5.73);
\node[anchor=east,font=\scriptsize,text=black!50] at (2.60,-6.0) {source};
\node[anchor=east,font=\scriptsize,text=black!50] at (2.60,-6.5) {target};
\fill[black!7] (5,-6.63) rectangle (10,-5.82);
\draw[black!35] (6.25,-5.88)--(6.25,-5.76);
\draw[black!35] (8.75,-5.88)--(8.75,-5.76);
\draw[bar,muted] (2.8,-6.0)--(10,-6.0) (5,-6.5)--(10,-6.5);
\draw[bar,muted,densely dashed] (2.8,-6.5)--(5,-6.5);
\draw[muted,fill=white,line width=.5pt] (5,-6.62)--(5.12,-6.5)--(5,-6.38)--(4.88,-6.5)--cycle;
\draw[bar] (10,-6.5)--(10.8,-6.5);
\draw[aid,muted] (5,-6.68)--(10,-6.68);
\draw[aid] (2.8,-6.84)--(10,-6.84);
\draw[black!35,line width=.8pt] (9.910,-6.090)--(10.090,-5.910) (9.910,-5.910)--(10.090,-6.090);
\draw[black!35,line width=.8pt] (9.910,-6.770)--(10.090,-6.590) (9.910,-6.590)--(10.090,-6.770);
\draw[black,line width=.8pt] (9.910,-6.930)--(10.090,-6.750) (9.910,-6.750)--(10.090,-6.930);
\draw[black,line width=.8pt] (10.710,-6.590)--(10.890,-6.410) (10.710,-6.410)--(10.890,-6.590);
\draw[rounded corners=1pt,black!35,fill=white] (5.78,-6.64) rectangle (7.75,-6.36);
\fill[black!35] (5.95,-6.5) circle (.055);
\fill[black!35] (6.18,-6.5) circle (.055);
\fill[black!35] (6.41,-6.5) circle (.055);
\fill[black!35] (6.64,-6.5) circle (.055);
\fill[black!35] (6.87,-6.5) circle (.055);
\fill[black!35] (7.1,-6.5) circle (.055);
\fill[black!35] (7.33,-6.5) circle (.055);
\fill[black!35] (7.56,-6.5) circle (.055);
\draw[line width=.85pt] (12.870,-6.250)--(12.965,-6.345)--(13.150,-6.110);
\draw[line width=.85pt] (13.720,-6.250)--(13.815,-6.345)--(14.000,-6.110);
\draw[black,line width=.8pt] (14.595,-6.355)--(14.805,-6.145) (14.595,-6.145)--(14.805,-6.355);
% Withdrawal is one shared vertical cut, not an implication arrow.
\draw[black!60,dash pattern=on 3pt off 2pt,line width=.65pt] (10,3.82)--(10,-7.14);
\draw[-{Stealth[length=1.8mm]},line width=.6pt] (2.8,-7.35)--(12.4,-7.35);
\node[anchor=east] at (2.45,-7.35) {time};
\end{tikzpicture}
\endgroup

\caption{One handoff and four look-alikes. The top row is a handoff that passes all three checks. Realization is checked on the target's own dynamics. Lineage is checked where source and target are live together, with no gap. Autonomy is checked after the source and every transfer aid are withdrawn. The rows below are separate proved controls. A detached copy and a gap have the same predictive readout as the source at the endpoint, and both fail lineage. The last two rows each change one thing in the passing handoff. Deleting one published cell fails only realization. Withdrawing the target's retained recharge fails only autonomy. A dash marks a check we do not evaluate for that case. Schematic. The rows follow the operational version, in which none of the three checks implies another. On the proved instances, conditions (4) to (7) hold in simplified form.}
\label{fig:three}
\end{figure}

The reason endpoint similarity cannot decide the matter is short, and we give
it informally first.
\begin{enumerate}
\item[(P1)] A copy can be built on a different system so that its final
predictive state matches the source's.
\item[(P2)] The copy can satisfy the conditions on its own system and keep
itself going without help.
\item[(P3)] If there is no stage at which source and copy are live together,
no handoff carries the source's ongoing organization into the copy.
\item[(C)] So realization and autonomy can both hold, with matching
endpoints, while lineage fails. Any criterion that looks only at endpoints, or
only at the target, cannot tell this copy from a continuation.
\end{enumerate}
Section~\ref{sec:copy} turns this argument into a finite countermodel, in which
the conditions hold in the simplified form described there, and
shows two further ways to match endpoints without continuing: a gap in which
the stream stops and later resumes, and a re-encoding of one system that never
changes system at all.

Three further results complete the picture. First, continuation is a
relation, not a function. When one source continues symmetrically into two
targets, no rule that respects the symmetry picks exactly one of them as the
continuer, and ancestry alone does not settle which, if either, is the
original (Section~\ref{sec:branching}). Second, what a target inherits
depends on the route the transfer took. Two histories can end in the same
current state and still differ in what they predict, so the target's present
must carry enough to tell such histories apart. Keeping route memory is one way
to do this, and a flatness or confluence argument is a sufficient, not a
necessary, ground for dropping it. Holonomy can rule out a present that is fixed by the shared boundary alone
(Section~\ref{sec:route}). Third, we exhibit one designed model, with a finite
physical state space and unbounded execution histories, on which all eight
conditions of the architecture, as adopted at interfaces fixed in advance, are
computed jointly from a single execution. On this model a weaker form of
continuation holds, continuation by ancestry, for both branches of a fission
and for the output of a guarded fusion, each once it has completed its own
acquisition (Section~\ref{sec:ancestry}). Its lineage is predictive ancestry,
which is weaker than live overlap. This is
not a certificate for the full architecture, whose complete panel of checks we
do not verify.

What we establish, and what we do not, can be said once. Every result is a
statement about finite systems, some with unbounded histories, and every
result is checked in the Lean theorem prover \citep{deMouraUllrich2021} using
its core library and standard axioms only. The results are structural. They
do not say that any system is conscious, that consciousness can or cannot be
moved between substrates, or that a continuation is the same person or
subject. The architecture's postulate of Substrate Indifference remains a
postulate, and our results about moving structure between finite systems do
not address it. Nothing here shows that any measurement of a real system identifies these
structures. Any step from our results to felt continuity needs hypotheses that
we state and do not assume. Finally, no single model in this
paper has both strict live-overlap lineage and all eight conditions at the
grade the architecture specifies. Exhibiting one is an open problem.

Our contributions are four.
\begin{enumerate}
\item We state realization, lineage and autonomy as separate predicates on
finite systems, and prove that a certificate checking them is sound
(Sections~\ref{sec:checks} and~\ref{sec:continuation}).
\item We give finite countermodels for the shortcuts: a detached copy, a gap,
a public shadow, agreement checked only between neighbouring pieces, a clock that ticks without
progress, a summed defect score, and success up to any finite horizon
(Sections~\ref{sec:copy} and~\ref{sec:checks}).
\item We show that continuation is relation-valued under fission, that no
symmetry-respecting rule selects a unique continuer, and that ancestry alone
underdetermines identity (Section~\ref{sec:branching}).
\item We exhibit a model on which all eight conditions are computed jointly,
and prove continuation by ancestry on it, including guarded fusion
(Section~\ref{sec:ancestry}).
\end{enumerate}

Section~\ref{sec:setting} fixes the setting and says how to read our claims.
Section~\ref{sec:copy} works through the copy example.
Section~\ref{sec:checks} gives the three checks, and
Section~\ref{sec:continuation} the continuation theorem.
Sections~\ref{sec:route} and~\ref{sec:branching} treat routes and branching,
Section~\ref{sec:ancestry} continuation by ancestry, and
Section~\ref{sec:modes} direct succession and the reach of the results.
Section~\ref{sec:objections} answers objections,
Section~\ref{sec:related} places the work, and
Section~\ref{sec:conclusion} concludes.
